\section{Error compuesto en modelos del mundo, MBPO y Dreamer}

\subsection*{a) Rollout real y simulado}

Con la acción $u_t=0$ y el mismo estado inicial para ambos sistemas, las
recurrencias son $s_{t+1}=1.2s_t$ y
$\hat{s}_{t+1}=1.2\hat{s}_t+0.10$. Si
$d_t=\hat{s}_t-s_t$ denota el error con signo, entonces
\[
    d_{t+1}=1.2d_t+0.10,
    \qquad
    d_t=0.10\sum_{j=0}^{t-1}1.2^j
       =0.50(1.2^t-1).
\]
La suma se obtiene al expandir la recurrencia desde $d_0=0$. Como
$d_t\geq 0$, el error absoluto $e_t$ coincide con $d_t$.

\begin{center}
    \begin{tabular}{cccc}
        \toprule
        Paso $t$ & Estado real $s_t$ & Estado simulado $\hat{s}_t$
                 & Error $e_t$ \\
        \midrule
        1 & $1.200$ & $1.300$ & $0.100$ \\
        2 & $1.440$ & $1.660$ & $0.220$ \\
        3 & $1.728$ & $2.092$ & $0.364$ \\
        \bottomrule
    \end{tabular}
\end{center}

El sesgo de cada transición se agrega al error previo multiplicado por
$1.2$. Por ejemplo, $e_3=1.2(0.220)+0.10=0.364$.

\subsection*{b) Rollouts cortos y ensambles en MBPO}

MBPO toma estados iniciales de la memoria de experiencias reales. Desde
cada uno ejecuta un rollout corto con acciones de la política actual y
transiciones generadas por el modelo. Estos fragmentos complementan la
experiencia real que recoge el agente y proporcionan datos para
actualizar la política. Iniciar cada fragmento en un estado observado
evita heredar el error de un rollout anterior. Limitar su longitud
reduce el número de pasos en que las predicciones propias se
retroalimentan y pueden llevar a regiones poco representadas en los
datos~\cite[secc.~5]{janner2019}.

MBPO entrena un ensamble de modelos probabilísticos sobre remuestras de
las transiciones reales. Cada miembro predice una distribución para el
estado siguiente y la recompensa. Al generar cada transición se elige
un miembro al azar, por lo que las muestras reflejan tanto la dispersión
de cada modelo como el desacuerdo entre modelos. Esto reduce la
dependencia de una sola predicción, aunque un sesgo compartido por el
ensamble puede persistir~\cite[secc.~5]{janner2019}.

\subsection*{c) Predicción latente en el RSSM de Dreamer}

El estado determinista $h_t=f_\phi(h_{t-1},z_{t-1},a_{t-1})$ resume el
historial y condiciona dos distribuciones sobre la misma variable
latente estocástica. Al entrenar el modelo del mundo con secuencias
reales, el posterior $q_\phi(z_t\mid h_t,o_t)$ incorpora la observación
$o_t$. El prior predictivo $p_\phi(\hat z_t\mid h_t)$ estima el latente
antes de recibir $o_t$. El término
\[
    D_{\mathrm{KL}}\bigl(
        q_\phi(z_t\mid h_t,o_t)\,\|\,p_\phi(\hat z_t\mid h_t)
    \bigr)
\]
acerca las predicciones del prior a las representaciones inferidas con
píxeles reales y favorece un espacio latente predecible a partir del
historial~\cite[secc.~4]{hafner2020}.

El Actor-Crítico parte de estados latentes inferidos de experiencias
reales. En los pasos imaginados, el actor elige acciones usando el
estado latente, la recurrencia calcula el siguiente $h_t$ y el prior
genera $\hat z_t$. El modelo de recompensa predice la recompensa a
partir de $(h_t,\hat z_t)$, y el crítico aprende a estimar retornos con
esas recompensas. La divergencia KL conecta las representaciones
aprendidas con el posterior y las que se generan con el prior. La
reconstrucción de imágenes puede aportar una señal al entrenar el
modelo del mundo. Durante los pasos imaginados del Actor-Crítico
bastan la transición y la predicción de recompensa en el espacio
latente~\cite[seccs.~3--4]{hafner2020}.
